% !TEX root = ../main.tex
%
% Methods subsection -- Training Protocol

\subsection{Training Protocol}
\label{subsec:training}

\textbf{Optimization.} The WSI-based models (M1--M4) are trained with AdamW
(learning rate $1\times10^{-5}$, weight decay $1\times10^{-1}$) for 30 epochs,
using a linear warmup over the first 3 epochs followed by cosine learning-rate
decay. The WSI-free baselines (M5--M7) contain only small MLP encoders without a
self-attention stack, and are trained with a higher learning rate
($1\times10^{-3}$, weight decay $1\times10^{-2}$) under the same warmup/cosine
schedule, with early stopping (patience of 20 epochs without validation C-index
improvement). Training uses bfloat16 mixed precision, which requires no loss
scaling since bfloat16 shares the same exponent range as float32.

\textbf{Batching and the survival loss.} Because the Cox partial-likelihood loss
requires comparing a patient against a risk set of other patients, training
batches are constructed at the patient level: gradients for all slides belonging
to one patient are accumulated before the model moves to the next patient, and
risk sets are formed from batches of 16 patients before each optimizer step.

\textbf{Regularization applied only during training.} Patch subsampling
($\kappa = 0.8$) and real-time tile augmentation are applied exclusively to the
training split; validation, internal test, and external test are always
evaluated on the full, unaugmented patch set. Model selection uses the
checkpoint with the highest validation C-index.

\textbf{Repeated runs.} Because random weight initialization alone can shift
external C-index by several hundredths (comparable to or larger than
differences between competing model variants), no model is evaluated from a
single training run. Each model is trained independently across multiple random
seeds (42, 84, 126) and, for external validation, in both training directions
(TCGA-PAAD $\rightarrow$ CPTAC-PDAC and CPTAC-PDAC $\rightarrow$ TCGA-PAAD),
yielding six runs per model for external comparisons; results are reported as
mean $\pm$ SD across runs. For internal evaluation, a stratified 5-fold
cross-validation protocol is additionally used, in which every patient is held
out as test exactly once and internal performance is computed by pooling
predictions across all five held-out folds.
